\clearpage\begingroup\SourceLeaf{337}\ArithmeticHead{6}{L}{STÄCKEL, GAUSS ALS GEOMETER.}
\renewcommand{\thefootnote}{\arabic{footnote})}
\begin{center}
\textls[100]{Abschnitt I.}\par\medskip
\textbf{Die Grundlagen der Geometrie.}\par\medskip
\textbf{2.}\par\medskip
\textls[80]{Allgemeines über die Arbeitsweise von \textsc{Gauss}.}
\end{center}
\noindent\hspace*{\parindent}Bei den Grundlagen der Geometrie zeigt sich in hohem Masse eine Erscheinung,
der wir bei \textsc{Gauss} wiederholt begegnen: der Reichtum der Gedanken,
die ihm, besonders in der Jugend, in solcher Fülle zuströmten, dass er ihrer
kaum Herr werden konnte (\textsc{Sartorius}, S. 78), steht in Gegensatz zu dem geringen
Umfang dessen, was er aufgezeichnet, ausgearbeitet und veröffentlicht
hat. Wenn daher auch der folgende Bericht über die Arbeitsweise von \textsc{Gauss}
mehr in eine (noch fehlende) Schilderung seiner gesamten wissenschaftlichen
Persönlichkeit als in eine Darlegung seiner Arbeiten auf einem Teilgebiet der
Mathematik zu gehören scheint, so dürfte er doch als Grundlage für das Verständnis
der folgenden Ausführungen nützlich sein, zumal dabei der Zusammenhang
mit den Grundlagen der Geometrie nicht aus dem Auge verloren wird.
\par
Ein erster Grund für die Erscheinung, auf die wir hingewiesen haben,
liegt darin, dass die Grösse des mathematischen Genies, das sich in \textsc{Gauss}
offenbarte, in der Vereinigung schöpferischer und kritischer Kraft wurzelt.
Diese Eigentümlichkeit erkennt man schon in der Dissertation, und sie zeigt
sich nicht weniger in den \textit{Disquisitiones arithmeticae}. In den späteren Veröffentlichungen
tritt die Kritik an den Leistungen anderer zurück, aber es bleibt
als auszeichnendes Merkmal die »\textsc{Gauss}sche Strenge«.
\par
Die \textsc{Gauss}sche Strenge erkennen wir schon äusserlich in der Form der
Darstellung. »Es war zu aller Zeit \textsc{Gauss}' Streben, seinen Untersuchungen
die Form vollendeter Kunstwerke zu geben; eher ruhete er nicht, und er hat
daher nie eine Arbeit veröffentlicht, bevor sie diese von ihm gewünschte,
durchaus vollendete Form erhalten hatte. Man dürfe einem Bauwerke, pflegte
er zu sagen, nach seiner Vollendung nicht’ mehr das Gerüste ansehen« (\textsc{Sartorius},
S. 82). Dieser Grundsatz spricht sich auch in dem Siegel aus, das
\textsc{Gauss} benutzte; es zeigt einen Baum mit wenigen Früchten und der Umschrift:
Pauca, sed matura.
\clearpage\SourceLeaf{338}\ArithmeticHead{7}{R}{DIE GRUNDLAGEN DER GEOMETRIE.}\setcounter{footnote}{0}
\noindent\hspace*{\parindent}In Briefen an \textsc{Schumacher}, \textsc{Encke} und \textsc{Bessel} hat \textsc{Gauss} sich darüber
geäussert, warum er von dieser klassischen Darstellungsart nicht abgehen
wollte und konnte.
\par
Als er nach Abschluss der geodätischen Messungen im Felde im Winter
1825/26 seine theoretischen Arbeiten wieder aufnimmt, klagt er am 21. November
1825 \textsc{Schumacher} gegenüber: »Der Wunsch, den ich immer bei meinen
Arbeiten gehabt habe, ihnen eine solche Vollendung zu geben, ut nihil amplius
desiderari possit\footnote{Diese Wendung findet sich bei \textsc{Euler}, siehe z. B. Nova acta acad. sc. Petrop. \textbf{4} (1786), 1789, S. 73.}, erschwert sie mir freilich ausserordentlich« (W. VIII, S. 400).
\textsc{Schumacher} antwortet am 2. Dezember 1825: »In Bezug auf Ihre Arbeiten und
den Grundsatz, ut nihil amplius desiderari possit, möchte ich fast wünschen,
und zum Besten der Wissenschaft wünschen, Sie hielten nicht so strenge daran.
Von dem unendlichen Reichtum Ihrer Ideen würde dann mehr uns werden
als jetzt, und mir scheint die Materie wichtiger als die möglich vollendetste
Form, deren die Materie fähig ist. Doch schreibe ich meine Meinung mit
Scheu hin, da Sie gewiss längst das pro und contra möglichst erwogen haben«
(Br. G.-\textsc{Sch.} II, S. 41). \textsc{Gauss} erwidert am 12. Februar 1826: »Ich war etwas
verwundert über Ihre Äusserung, als ob mein Fehler darin bestehe, die Materie
zu sehr der vollendeten Form hintanzusetzen. Ich habe während meines
ganzen wissenschaftlichen Lebens immer das Gefühl gerade vom Gegenteil
gehabt, d. i. ich fühle, dass oft die Form vollendeter hätte sein können und
dass darin Nachlässigkeiten zurückgeblieben sind. Denn so werden Sie es
doch nicht verstehen, als ob ich \textls[100]{mehr für die Wissenschaft leisten
würde}, wenn ich mich damit begnügte, einzelne Mauersteine, Ziegel etc. zu
liefern, anstatt eines Gebäudes, sei es nun ein Tempel oder eine Hütte, da
gewissermassen das Gebäude auch nur Form der Backsteine ist. Aber ungern
stelle ich ein Gebäude auf, worin Hauptteile fehlen, wenngleich ich wenig
auf den äusseren Aufputz gebe. Auf keinen Fall aber, wenn Sie sonst mit
Ihrem Vorwurf auch Recht hätten, passt er auf meine Klagen über die gegenwärtigen
Arbeiten, wo es nur das gilt, was ich Materie nenne; und ebenso
kann ich Ihnen bestimmt versichern, dass, wenn ich gern auch eine gefällige
Form gebe, \textls[100]{diese} vergleichungsweise nur sehr wenig Zeit und Kraft in Anspruch
nimmt oder bei früheren Arbeiten genommen hat« (Br. G.-\textsc{Sch.} II, S. 46).
\clearpage\SourceLeaf{339}\ArithmeticHead{8}{L}{STÄCKEL, GAUSS ALS GEOMETER.}
\noindent\hspace*{\parindent}Als \textsc{Gauss} bald darauf an \textsc{Schumacher} eine kleine Abhandlung über den
Heliotropen für die Astronomischen Nachrichten sendet (W. IX, S. 472), fügt
er hinzu: »Diesmal habe ich gewiss den Vorwurf nicht verdient, als ob ich der
Form auf Kosten der Materie \textls[100]{zuviel} eingeräumt hätte, sondern eher das Gegenteil«
(Brief vom 28. November 1826, Br. G.-\textsc{Sch.} II, S. 81), und \textsc{Schumacher}
sieht sich jetzt veranlasst, seine Meinung ausführlich auseinanderzusetzen. Am
Schluss heisst es: »Ich glaubte, dies Ausfeilen könne ebenso gut ein anderer
tun, und darin \textls[100]{kann} ich mich geirrt haben; worin ich mich aber nicht geirrt
habe, ist die Behauptung, dass Sie das Erfinden nicht einem andern übertragen
können. Jedes Jahr Ihres Lebens mehrt die Ihnen nur verständlichen Andeutungen
neuer Ideen. Soll alles dieses verloren sein?« (Brief vom 2. Dezember
1826, Br. G.-\textsc{Sch.} II, S. 83).
\par
\textsc{Gauss} verhielt sich solchen Anregungen gegenüber durchaus ablehnend.
»Ich weiss«, schreibt er am 18. August 1832 an \textsc{Encke}, »dass einige meiner
Freunde wünschen, dass ich weniger in diesem Geiste arbeiten möchte: das
wird aber nie geschehen; ich kann einmal an Lückenhaftem keine rechte
Freude haben, und eine Arbeit, an der ich keine Freude habe, ist mir nur
eine Qual. Möge auch jeder in \textls[100]{dem} Geiste arbeiten, der ihm am meisten
zusagt« (W. XI\textsubscript{1}, S. 84).
\par
Am 15. Januar 1827 berichtet er seinem Freunde \textsc{Schumacher}, er sei mit
der Ausarbeitung der Abhandlung über die krummen Flächen ein gut Stück
vorgerückt. »Ich finde dabei viele Schwierigkeiten, allein das, was man \textls[100]{Ausfeilen}
oder \textls[100]{Form} mit Recht nennen könnte, ist doch keineswegs, was erheblich
aufhält (wenn ich die Sprödigkeit der lateinischen Sprache ausnehme),
vielmehr ist es die innige Verkettung der Wahrheiten in ihrem Zusammenhange,
und eine solche Arbeit ist erst dann gelungen, wenn der Leser die
grosse Mühe, die bei der Ausführung stattgefunden hat, gar nicht mehr erkennt.
Ich kann daher nicht leugnen, dass ich keinen recht klaren Begriff
davon habe, wie ich meine Arbeiten solcher Art anders, als ich gewohnt bin,
ausführen könnte, ohne, wie ich mich schon einmal ausgedrückt habe, Mauersteine
anstatt eines Gebäudes zu liefern. Ich habe wohl zuweilen versucht,
über diesen oder jenen Gegenstand bloss Andeutungen ins Publikum zu
bringen; entweder aber sind sie von Niemand beachtet oder wie z. B. einige
Äusserungen in einer Rezension G. G. Anz. 1816, p. 619 [W. IV, S. 364,
\clearpage\SourceLeaf{340}\ArithmeticHead{9}{R}{DIE GRUNDLAGEN DER GEOMETRIE.}\HistoricalFoot{X\textsubscript{2} Abh. 4.}{2}
\noindent VIII, S. 170], es ist mit Kot danach geworfen. Also, insofern von wichtigen
Gegenständen die Rede ist, etwas im Wesen Vollendetes oder gar nichts«
(Br. G.-\textsc{Sch.} II, S. 93). Solche Andeutungen finden sich zahlreich in den
Jugendwerken, sie fehlen aber auch nicht in den späteren Schriften. Wie
sorgfältig \textsc{Gauss} dabei verfuhr, zeigt der Brief an \textsc{Encke} vom 18. August 1832,
wo es heisst: »Es ist von jeher mein gewissenhaft befolgter Grundsatz gewesen,
solche Andeutungen, die aufmerksame Leser in jeder meiner Schriften in grosser
Menge finden (sehen Sie z. B. meine \textit{Disquis. arithmet.} pag. 593 [art. 335]) stets
dann erst zu machen, wenn ich den Gegenstand für mich selbst ganz abgemacht
habe« (W. XI\textsubscript{1}, S. 84). Hiernach wird man im besonderen die vorher
erwähnten Andeutungen in den Göttinger Anzeigen vom Jahre 1816 zu bewerten
haben, die sich auf die Unbeweisbarkeit des Parallelenaxioms beziehen.
\par
Es ist ein merkwürdiger Zufall, dass \textsc{Gauss} bald, nachdem er sich bei
\textsc{Schumacher} über die Erfolglosigkeit seiner Andeutungen beklagt hatte, am
24. Juli und 14. August 1827 (Br. G.-\textsc{Sch.} II, S. 105, 111) durch seinen Freund
die beiden Briefe \textsc{Jacobis} erhielt, mit denen dessen Untersuchungen über die
elliptischen Funktionen beginnen (\textsc{Jacobi}, Werke I, S. 29), und dass er nicht
lange danach \textsc{Abels} \textit{Recherches} kennen lernte, die ihm von seinen eigenen
Untersuchungen »wohl ein Drittel vorwegnahmen« (Brief vom 30. Mai 1828,
Br. G.-\textsc{Sch.} II, S. 177). Der ausschlaggebende Einfluss, den die berühmte
Stelle im art. 335 der \textit{Disquisitiones arithmeticae} (W. I, S. 412) auf \textsc{Abel} und
\textsc{Jacobi} geübt hat, ist anerkannt. Hier hat ein von \textsc{Gauss} ausgestreutes Samenkorn
hundertfältige Frucht getragen, und auch andere Andeutungen sind nicht
auf steinigen Boden gefallen.
\par
Fast ein Vierteljahrhundert später ist derselbe Streitpunkt zwischen den
beiden Freunden noch einmal aufgetaucht, als nämlich \textsc{Schumacher} in den
Astronomischen Nachrichten \textsc{Jacobis} Bearbeitung der \textsc{Carlini}schen Abhandlung
über die \textsc{Kepler}sche Gleichung abdruckte und \textsc{Gauss} jenem mitteilte (Brief
vom 4. Dezember 1849, Br. G.-\textsc{Sch.} VI, S. 51), er habe die Aufgabe schon
vor langer Zeit »auf eine ohne allen Vergleich kürzere Art aufgelöst« (W. X\textsubscript{1},
S. 420—428). »Wenn ich nicht wüsste«, hatte \textsc{Schumacher} geantwortet, »wieviel
Zeit Ihnen die letzte Feile Ihrer Arbeiten kostet, so würde ich um Ihre
Abhandlung bitten« (Brief vom 5. Dezember 1849, Br. G.-\textsc{Sch.} VI, S. 52).
\textsc{Gauss} erwidert, er sei nicht abgeneigt, eine ihm zu Teil werdende Musse
\clearpage\SourceLeaf{341}\ArithmeticHead{10}{L}{STÄCKEL, GAUSS ALS GEOMETER.}\setcounter{footnote}{0}
\noindent zur Ausarbeitung einer Abhandlung über den Gegenstand zu verwenden; es
werde aber erhebliche Zeit erfordert werden, um die ganze Theorie in einer
ihm selbst genügenden Gestalt auszuführen. »Sie sind ganz im Irrtum, wenn
Sie glauben, dass ich darunter nur die letzte Politur in Beziehung auf Sprache
und Eleganz der Darstellung verstehe. Diese kosten vergleichungsweise nur
unbedeutenden Zeitaufwand; was ich meine, ist die \textls[100]{innere} Vollkommenheit.
In manchen meiner Arbeiten sind solche Inzidenzpunkte, die mich jahrelanges
Nachdenken gekostet haben, und deren in kleinem Raum konzentrierte Darstellung
nachher niemand die Schwierigkeit anmerkt, die erst überwunden
werden muss[te]« (Brief vom 5. Februar 1850, Br. G.-\textsc{Sch.} VI, S. 58).
\par
Ähnliche Äusserungen finden sich in dem Briefe an \textsc{Bessel} vom 28. Februar
1839 (Br. G.-\textsc{Bessel}, S. 524); ihnen gegenüber vertritt \textsc{Bessel} in dem
Briefe vom 28. Juni 1839 (Br. G.-\textsc{Bessel}, S. 526) mit grosser Wärme den
Standpunkt, den \textsc{Schumacher} in dem Briefe vom 2. Dez. 1826 eingenommen
hatte.
\par
Die vollendete Darstellung, bei der \textsc{Archimedes} und \textsc{Newton} für \textsc{Gauss}
die Vorbilder waren, sollte nur das äussere Zeichen der inneren Vollkommenheit
sein, und hier erst gewinnt das Wort von der \textsc{Gauss}schen Strenge seine
wahre Bedeutung. Von den Geometern des 18. Jahrhunderts war in der
Freude über die Fülle neuer Entdeckungen, zu denen die Infinitesimalrechnung
die Mittel bot, die Sicherung der Grundlagen ausser Acht gelassen
worden. Sehr stark tritt das bei \textsc{Euler} hervor, bei dem gerade die grundlegenden
Betrachtungen viel zu wünschen übrig lassen\footnote{Vgl. etwa L. \textsc{Schlesinger} und F. \textsc{Engel} in der Vorrede zu \textsc{Eulers} \textit{Institutiones calculi integralis}, Opera omnia, ser. I, vol. 11, Leipzig 1913, S. XIII.}.
Dagegen finden sich
schon bei \textsc{d'Alembert} Ansätze zu einer kritischen oder besser skeptischen
Auffassung, und \textsc{Lagrange} hat in der \textit{Théorie des fonctions analytiques} geradezu
das Ziel erstrebt, den Beweisen den Charakter einleuchtender Gewissheit und
Strenge zu geben, der die Lösungen der Alten auszeichnet\footnote{J. L. \textsc{Lagrange}, \textit{Théorie des fonctions analytiques}, Paris 1797; Oeuvres, t. 9, S. 184.}.
Der »rigor apud
veteres consuetus« ist es, den der junge \textsc{Gauss} im bewussten Gegensatz zu
den Gepflogenheiten des 18. Jahrhunderts auf seine Fahne geschrieben hat\footnote{C. F. \textsc{Gauss}, \textit{Disquisitiones arithmeticae}, Lipsiae 1801, Praefatio; W. I, S. 5.}.
Im hohen Alter hat er \textsc{Schumacher} gegenüber seine Überzeugung mit folgenden
\clearpage\SourceLeaf{342}\ArithmeticHead{11}{R}{DIE GRUNDLAGEN DER GEOMETRIE.}\HistoricalFoot{}{2\textsuperscript{*}}
\noindent Worten ausgesprochen: »Es ist der Charakter der Mathematik der neueren
Zeit (im Gegensatz gegen das Altertum), dass durch unsere Zeichensprache
und Namengebungen wir einen Hebel besitzen, wodurch die verwickeltsten
Argumentationen auf einen gewissen Mechanismus reduziert werden. An
Reichtum hat dadurch die Wissenschaft unendlich gewonnen, an Schönheit
und Solidität aber, wie das Geschäft gewöhnlich betrieben wird, eben so sehr
verloren. Wie oft wird jener Hebel eben nur mechanisch angewandt, obgleich
die Befugnis dazu in den meisten Fällen gewisse stillschweigende Voraussetzungen
impliziert. Ich fordere, man soll bei allem Gebrauch des Kalküls,
bei allen Begriffsverwendungen sich immer der ursprünglichen Bedingungen
bewusst bleiben, und \textls[100]{alle} Produkte des Mechanismus niemals über die klare
Befugnis hinaus als Eigentum betrachten. Der gewöhnliche Gang ist aber
der, dass man für die Analysis einen Charakter der Allgemeinheit in Anspruch
nimmt und dem \textls[100]{Andern}, der so herausgebrachte Resultate noch nicht für
bewiesen anerkennt, zumutet, er solle das Gegenteil nachweisen. Die Zumutung
darf man aber nur an den stellen, der seinerseits \textls[100]{behauptet}, ein
Resultat sei falsch, nicht aber dem, der ein Resultat nicht für bewiesen anerkennt,
welches auf einem Mechanismus beruhet, dessen ursprüngliche, wesentliche
Bedingungen in dem vorliegenden Fall gar nicht zutreffen« (Brief an
\textsc{Schumacher} vom 1. September 1850, W. X\textsubscript{1}, S. 434).
\par
Ein zweiter Grund für das Missverhältnis zwischen dem Reichtum an
Gedanken, die »bei der unglaublichen Produktivität in dem mächtigen Gehirn
auftauchten« (\textsc{Sartorius}, S. 79), und dem verhältnismässig geringen Umfang
der rein mathematischen Veröffentlichungen von \textsc{Gauss} liegt in Hemmungen
innerer und äusserer Art, die bei seiner Art des Arbeitens dem Druckfertigmachen
entgegenstanden.
\par
In dem schon erwähnten Briefe an \textsc{Bessel} vom 28. Februar 1839 hatte
\textsc{Gauss} mit einer bei ihm ungewöhnlichen Heftigkeit des Tones hervorgehoben,
er brauche zum Ausarbeiten »Zeit, viel Zeit, viel mehr Zeit, als Sie sich wohl
vorstellen mögen. Und meine Zeit ist vielfach beschränkt, sehr beschränkt«.
Solche Klagen über Mangel an Zeit für die theoretischen Untersuchungen
wiederholen sich beständig in den Briefen. Die glücklichste Zeit seines Lebens
sind wohl jene neun Jahre von 1799 bis 1807 gewesen, die er als Schützling
des »edlen Fürsten, dem er alles, was er war, verdankte« (Brief an \textsc{Olbers}
\clearpage\SourceLeaf{343}\ArithmeticHead{12}{L}{STÄCKEL, GAUSS ALS GEOMETER.}\setcounter{footnote}{0}
\noindent vom 23. Februar 1802, Br. G.-O. 1, S. 14) in Braunschweig zugebracht hat.
Noch im Alter hat er dieser Jahre mit Rührung und Dankbarkeit gedacht.
So schreibt er am 15. Februar 1845 an \textsc{Encke} über \textsc{Eisenstein}, der damals
mit Unterstützung des Königs von Preussen in freier Musse seinen mathematischen
Forschungen nachging: »Er lebt noch in der glücklichen Zeit, wo
er sich ganz seiner Begabung hingeben kann, ohne dass er nötig hätte, sich
durch irgend etwas Fremdartiges stören zu lassen. Ich werde lebhaft an die
— längst verflossenen — Jahre erinnert, wo ich in ähnlichen Verhältnissen
lebte. Von der andern Seite erfordern auch gerade die rein mathematischen
Spekulationen eine unverkümmerte und unzerstückelte Zeit« (Brief im \textsc{Gauss}-Archiv).
\par
Die Pflichten der Professur haben schwer auf \textsc{Gauss} gelastet, zunächst
sein Amt als Leiter der Göttinger Sternwarte. »So sehr ich die Astronomie
liebe«, schreibt er am 28. Juni 1820 an \textsc{Bessel} (Br. G.-\textsc{Bessel}, S. 353), »fühle ich
doch das Beschwerliche des Lebens eines praktischen Astronomen, ohne Hilfe,
oft nur zu sehr, am peinlichsten aber darin, dass ich darüber fast gar nicht
zu irgend einer zusammenhängenden grösseren theoretischen Arbeit kommen
kann«.
\par
Hierzu traten seit 1821 die geodätischen Messungen, und wenn auch die
mühsamen und zeitraubenden Arbeiten im Felde für \textsc{Gauss} selbst mit dem
Jahre 1825 beendet waren, so behielt er doch die Oberaufsicht über die Triangulationen
und führte die abschliessenden Rechnungen. »Mehr als zwanzig
Jahre hindurch«, sagt \textsc{Gaede}\footnote{\textsc{Gaede}, \textit{Beiträge zur Kenntnis von Gauss' praktisch-geodätischen Arbeiten}, Zeitschrift für Vermessungswesen, Bd. 14, 1885; auch als selbständiges Werk, Karlsruhe 1885, erschienen, S. 68.}, »hat \textsc{Gauss} unter der ermüdenden Last dieses
Geschäftes gelebt und gelitten, welches, wenn einmal in Gang gebracht und
in zweckmässiger Weise schematisch organisiert, von jedem andern ebenso
gut hätte besorgt werden können, während \textsc{Gauss} durch die massenhafte, und
sobald die Methode feststand, im Wesentlichen nur noch mechanische Rechenarbeit
der Musse verlustig ging, deren er für seine schöpferische Tätigkeit
auf spekulativem Gebiet, nach seinem eigenen Zeugnis, in hohem Masse bedurfte«.
\par
Dazu kam die Verpflichtung, Vorlesungen zu halten. »Für eine mathematische
Lehrstelle hat er eine ganz entschiedene Abneigung«, hatte \textsc{Olbers}
\clearpage\SourceLeaf{344}\ArithmeticHead{13}{R}{DIE GRUNDLAGEN DER GEOMETRIE.}
\noindent am 3. November 1802 an \textsc{Heeren} in Göttingen geschrieben, als es sich um
eine Berufung von \textsc{Gauss} an die dortige Universität handelte, »sein Lieblingswunsch
ist, Astronom bei irgend einer Sternwarte zu werden, um seine ganze
Zeit zwischen Beobachtungen und feinen, tiefsinnigen Untersuchungen zur Erweiterung
der Wissenschaft teilen zu können« (\textsc{Sartorius}, S. 31). Allein seine
Stellung an der Universität brachte es mit sich, dass er »das Handwerk eines
Professors« (\textsc{Sartorius}, S. 96) ausüben musste. Er hat es mit der ihm eigenen
Gewissenhaftigkeit getan, aber schon in dem Briefe an \textsc{Bessel} vom 27. Januar
1816 (Br. G.-\textsc{Bessel}, S. 232) nennt er das Kollegienlesen »ein sehr lästiges,
undankbares Geschäft«, und ganz besonders bitter werden seine Klagen, als
die Last der geodätischen Messungen hinzukommt. Die in Aussicht stehende
Berufung nach Berlin veranlasst ihn 1824 zu dem Ausruf: »Ich bin ja hier
so weit davon entfernt, Herr meiner Zeit zu sein. Ich muss sie teilen zwischen
Kollegia lesen (wogegen ich von jeher einen Widerwillen gehabt habe, der,
wenn auch nicht entstanden, doch vergrössert ist durch das Gefühl, welches
mich immer dabei begleitet, meine Zeit wegzuwerfen) und praktisch astronomische
Arbeiten.\ .\ .\ .\quad Was bleibt mir also für solche Arbeiten, auf die ich
selbst einen höhern Wert legen könnte, als flüchtige \textls[100]{Nebenstunden}? Ein
anderer Charakter als der meinige, weniger empfindlich für unangenehme Eindrücke,
oder ich selbst, wenn manches andere anders wäre, als es ist, würde
vielleicht auch solchen Nebenstunden noch mehr abgewinnen, als ich es im
\textls[100]{allgemeinen} kann« (Brief an \textsc{Bessel} vom 14. März 1824, Br. G.-\textsc{Bessel},
S. 428). Es liessen sich den Briefen an die vertrauten Freunde noch zahlreiche
Klagen dieser Art entnehmen. Hier möge nur noch eine Stelle aus
dem Briefe an \textsc{Olbers} vom 19. Februar 1826 (Br. G.-O. 2, S. 438) angeführt
werden: »\textls[100]{Unabhängigkeit}, das ist das grosse Losungswort für die Geistesarbeiten
in die Tiefe. Aber wenn ich meinen Kopf voll von in der Luft
schwebenden geistigen Bildern habe, die Stunde heranrückt, wo ich Kollegien
lesen muss, so kann ich Ihnen nicht beschreiben, wie angreifend das Abspringen,
das Anfrischen heterogener Ideen für mich ist, und wie schwer mir
oft Dinge werden, die ich unter andern Umständen für eine erbärmliche ABC-Arbeit
halten würde.\ .\ .\ .\quad Inzwischen, lieber \textsc{Olbers}, will ich Sie nicht mit
Klagen über Dinge [er]müden, die nicht zu ändern sind; meine ganze Stellung
\clearpage\endgroup
